% chapters/ch3_study1.tex — 第 3 章 研究内容一
%
% 对应小论文：MMM2027（RoGS-SLAM，camera-ready 14 页）
% 改写纪律：不直接翻译会议论文；claim 强度按大论文展开；实验数字须从
% results/evidence 工件重取（见 ../WRITING_RULES.md 六、七）。
% 本次只写 引言 + 方法；实验小节待数字核对后补，小结留骨架。

\chapter{动态场景下的可靠性引导观测处理}
\label{ch:study1}

\section{引言}

同步定位与地图构建（simultaneous localization and mapping，SLAM）在动态场景里
遇到的困难不在静态几何，而在观测本身不可信。行人走进视野，它在图像上的位移既可能是自身
运动，也可能是相机运动造成的视差；把两者混在一起交给跟踪器，位姿估计会被拖偏，
把两者一起当作噪声剔掉，又可能连带着删掉真实结构。动态 SLAM 的第一道关口因此不是
怎样建图，而是怎样判断一次观测该信多少。

早期做法沿着「先分割、后剔除」的路线展开。DynaSLAM~\cite{DynaSLAM} 用语义分割与
多实例跟踪圈出运动像素，直接从数据关联中去掉，再对被遮挡的背景做补绘。这条路线
简单有效，代价是把观测当成二值开关，一旦掩码出错，真实结构也一并被抹掉。语义分割
在两类场景下格外脆弱，一是物体运动但语义上不属于「动态类」的情形，二是分割模型
没有见过的场景，这两类失败在野外数据里都不罕见。

另一种思路是保留观测，只调整它的权重。既然错误的观测会持续给出与几何不一致的残差，
那么就用残差本身判断可信度，让可疑的观测平滑降权而不是骤然删除。研究内容一沿这条
思路，回答第 1 章提出的问题，即哪些观测可以进入跟踪和建图。本章的方法由四个组件
组成，合起来把「信多少」落成逐像素、逐帧的可计算量。

\section{方法}
\label{sec:study1-method}

\subsection{可靠性信号}

可靠性信号的核心是构造两个可以互相校验的流量。观测流来自图像上相邻关键帧之间的
像素位移，由光流给出；刚体流则把上一帧的深度反投影到世界坐标，再用当前位姿投回
图像平面，得到如果场景完全静止时像素应该出现的位置。两者之差衡量的是「观测不符合
刚体假设」的程度。场景里的移动物体、未被建模的形变、深度噪声都会抬高这个残差，
被遮挡后重新露出的区域同样如此。

残差本身不能直接当权重用。行人在视野边缘时残差会短暂升高，人物静止不动时残差
接近零，单帧的读数波动很大。方法把逐像素残差转成 Cauchy 分布的似然，用来做只降不升
的重加权。Cauchy 的尾部比高斯厚，允许少数极大的残差被当作离群点压下去而不影响
大部分像素的权重；只降不升保证可靠性信号永远不会把本已可疑的观测抬得比原始观测更高，
避免信号自身的错误被放大成增益。

\subsection{鲁棒重加权}

跟踪环节的残差用 Huber 损失做迭代重加权。Huber 在残差较小时保持二次型，残差较大时
切换为线性，从而限制单次更新的影响力。这一步与可靠性信号互补，后者处理「像素是否
符合刚体假设」，前者处理「位姿更新是否被少数大残差带偏」。两者都作用在跟踪的入口，
不触碰建图端的高斯表示。

\subsection{关键帧覆盖}

关键帧的选取在人群场景里容易被运动物体带偏。行人经过时画面变化剧烈，按常规的图像
差异准则会连续插入关键帧，结果是关键帧密集堆在同一个视野，重建却缺少长程覆盖。
方法改用统一的覆盖准则，让关键帧分布服从场景几何而不是局部图像差异，同时限制相邻
关键帧的间隔上限，避免长距离漂移无人校正。

\subsection{lineage 生命周期}

第四个组件处理高斯的来去。运动物体在视野里停留的时段内，会有一批高斯被临时创建
来拟合它；物体离开后，这些高斯没有观测支撑，却可能因为不透明度尚未衰减而残留在
地图里。方法给每个高斯记录来源与代际，把「何时可以提交、何时应当回收」变成可审计的
生命周期，而不是依赖统一的不透明度阈值。语义掩码作为可选组件保留，与可靠性信号
叠加使用，不单独承担判断。

四个组件的共同点是把动态观测当作概率对象处理：不判断它是否动态，只估计它有多可信；
不一次性删除，而是持续降权并记录来源。第 3.3 节报告这些组件在标准动态序列上的
效果与边界，数字的工件来源见各表脚注。

% TODO(写作队列)：3.3 实验（数字须从 RoGS-SLAM/results/evidence 重取并与表逐格核对）；
% 3.4 小结（按 WRITING_RULES 四，避免复述式小结，写边界与适用条件）。
